\documentclass[12pt]{article}
% Préambule commun à tous les documents extraits de « La Revue CAPES »
\usepackage[T1]{fontenc}
\usepackage[utf8]{inputenc}
\usepackage{lmodern}
\usepackage{amsmath,amssymb,amsthm,amsfonts,mathrsfs}
\usepackage{setspace}
\usepackage{listings}
\usepackage{pgfplots}
\pgfplotsset{compat=1.18}
\usepackage{longtable}
\usepackage{adjustbox}
\usepackage{lettrine}
\usepackage{tkz-tab}
\usepackage{changepage}
\usepackage{pdflscape}
\usepackage{graphicx}
\usepackage{tikz}
\usepackage{xcolor}
\usepackage[a4paper,top=2cm,bottom=2.5cm,left=2.5cm,right=2cm]{geometry}
\usepackage{fancyhdr}
\usepackage{draftwatermark}
\SetWatermarkText{La RevueCapes}
\SetWatermarkLightness{0.9}
\SetWatermarkScale{2}
\usepackage{hyperref}
\hypersetup{colorlinks=true,linkcolor=blue!50!black,urlcolor=blue!50!black}

\definecolor{revue}{RGB}{20,60,120}

\pagestyle{fancy}
\fancyhf{}
\renewcommand{\headrulewidth}{0.4pt}
\setlength{\headheight}{15pt}

% \entete{Matière}{Titre}{Type de document}
\newcommand{\entete}[3]{%
  \fancyhead[L]{\small\textsc{La Revue CAPES} -- #1}%
  \fancyhead[R]{\small #2}%
  \fancyfoot[C]{\thepage}%
  \thispagestyle{plain}%
  \begin{center}
    {\color{revue}\rule{\textwidth}{1.2pt}}\\[4pt]
    {\small\textsc{Concours directs CAP-CEG / CAPES -- Burkina Faso}}\\[6pt]
    {\LARGE\bfseries\color{revue} #2}\\[6pt]
    {\large\itshape #1 -- #3}\\[2pt]
    {\color{revue}\rule{\textwidth}{1.2pt}}
  \end{center}
  \vspace{0.5cm}%
}

% Encadré pour les remarques ajoutées dans les corrigés
\newcommand{\remarque}[1]{\par\medskip\noindent\fcolorbox{revue}{revue!6}{\parbox{\dimexpr\linewidth-2\fboxsep-2\fboxrule}{\textbf{\color{revue}Remarque.} #1}}\par\medskip}


\begin{document}
\entete{Mathématiques}{Corrigé détaillé -- Session 2017}{Correction}
\begin{spacing}{1.5}

\textbf{Exercice 1}

On considère la matrice carrée d'ordre 3 définie par  : 

$M=\begin{pmatrix}
11 & -5 & 5 \\ 
-5 & 3 & -3 \\ 
5 & -3 & 3
\end{pmatrix}$

1. Déterminons le polynôme caractéristique de $M$, ainsi que son spectre.

\begin{align*}
P_M(x)=\det(M-xI_3)&=\begin{vmatrix}
11-x & -5 & 5 \\ 
-5 & 3-x & -3 \\ 
5 & -3 & 3-x
\end{vmatrix}c_3\leftarrow c_3+c_2\\
&=\begin{vmatrix}
11-x & -5 & 0 \\ 
-5 & 3-x & -x \\ 
5 & -3 & -x
\end{vmatrix}L_2\leftarrow L_2-L_3\\
&=x\begin{vmatrix}
11-x & -5 & 0 \\ 
-10 & 6-x & 0 \\ 
5 & -3 & -1
\end{vmatrix}\\
&=-x\begin{vmatrix}
11-x & -5 \\
-10 & 6-x 
\end{vmatrix}\\
&=-x((11-x)(6-x)-50)\\
&=-x(66-11x-6x+x^2-50)\\
&=-x(x^2-17x+16)\\
\end{align*}

Factorisation de $x^2-17x+16$

$\Delta=225=(15)^2$

$x=\frac{17-15}{2}=1$ ou $x=\frac{17+15}{2}=16$

Donc $P_M(x)=-x(x-1)(x-16)$.

Ainsi spectre de $M$ est $Sp(M)=\{0;1;16\}$

2. Les sous espaces propres associés aux différentes valeurs propres.

Soit $E_\lambda$ le sous-espace propre lié à la valeur propre $\lambda$

\[
E_\lambda=\{u=(x,y,z)\in \mathbb{R}^3/Mu=\lambda u\}
\]

\medskip
$\bullet$ Pour $\lambda=0$, on a :

$\forall u\in E_0,Mu=0\Rightarrow \begin{pmatrix}
11 & -5 & 5 \\ 
-5 & 3 & -3 \\ 
5 & -3 & 3
\end{pmatrix}\begin{pmatrix}
x\\
y\\
z
\end{pmatrix}=\begin{pmatrix}
0\\
0\\
0
\end{pmatrix}\Rightarrow \begin{cases}
11x-5y+5z=0\\
-5x+3y-3z=0\\
5x-3y+3z=0
\end{cases}\Rightarrow \begin{cases}
11x-5y+5z=0(1)\\
5x-3y+3z=0(2)
\end{cases}$

$(2)\Longrightarrow x=\frac{3}{5}(y-z)$

$(1)\Longrightarrow \frac{33}{5}(y-z)-5y+5z=0\Longrightarrow  \frac{8}{5}y-\frac{8}{5}=0\Longrightarrow y-z=0\Longrightarrow y=z$ $\Rightarrow x=0$

$\Rightarrow u=\begin{pmatrix}
0\\
y\\
y
\end{pmatrix}=y\begin{pmatrix}
0\\
1  \\
1
\end{pmatrix}\Rightarrow E_0=\langle \begin{pmatrix}
0\\
1\\
1
\end{pmatrix}\rangle$

$\bullet$ Pour $\lambda=1$, on a :

$\begin{pmatrix}
11 & -5 & 5 \\ 
-5 & 3 & -3 \\ 
5 & -3 & 3
\end{pmatrix}\begin{pmatrix}
x\\
y\\
z
\end{pmatrix}=\begin{pmatrix}
x\\
y\\
z
\end{pmatrix}=\begin{cases}
10x-5y+5z=0(1)\\
-5x+2y-3z=0(2)\\
5x-3y+2z=0(3)
\end{cases}$

$(2)+(3)\Rightarrow z=-y (4)$

$(4)$ dans $(1)\Longrightarrow x=y$

$\Rightarrow u=\begin{pmatrix}
y\\
y\\
-y
\end{pmatrix}=y\begin{pmatrix}
1\\
1\\
-1
\end{pmatrix}\Rightarrow E_{1}=\langle \begin{pmatrix}
1 \\
1 \\
-1
\end{pmatrix}\rangle$

$\bullet$ Pour $\lambda=16$, on a :

$\begin{pmatrix}
11 & -5 & 5 \\ 
-5 & 3 & -3 \\ 
5 & -3 & 3
\end{pmatrix}\begin{pmatrix}
x\\
y\\
z
\end{pmatrix}=16\begin{pmatrix}
x\\
y\\
z
\end{pmatrix}=\begin{cases}
-5x-5y+5z=0(1)\\
-5x-13y-3z=0(2)\\
5x-3y-13z=0(3)
\end{cases}$

$(2)+(3)\Rightarrow z=-y (4)$

$(4)$ dans $(1)\Longrightarrow x=-2y$

$\Rightarrow u=\begin{pmatrix}
-2y\\
y\\
-y
\end{pmatrix}=-y\begin{pmatrix}
2\\
-1\\
1
\end{pmatrix}\Rightarrow E_{16}=\langle \begin{pmatrix}
2 \\
-1 \\
1
\end{pmatrix}\rangle$


3. $M$ est diagonalisable car le polynôme caractéristique de $M$ est scindé en racines simples. Ou autrement car les valeurs propres sont distinctes deux à deux.

\medskip
\textbf{Diagonalisons $M$}

On a la base $B=\langle \begin{pmatrix}
0\\
1\\
1
\end{pmatrix},\begin{pmatrix}
1 \\
1 \\
-1
\end{pmatrix},\begin{pmatrix}
2 \\
-1 \\
1
\end{pmatrix}\rangle$.

Dans la base $B$, $D=\begin{pmatrix}
0 & 0 & 0 \\ 
0 & 1 & 0 \\ 
0 & 0 & 16
\end{pmatrix}$ et $P= \begin{pmatrix}
0 & 1  & 2 \\ 
1 & 1  & -1\\ 
1 & -1 & 1 
\end{pmatrix}$

$M=PDP^{-1}$

\medskip
Calculons $P^{-1}$

$P^{-1}=\frac{1}{det(P)}t_{com(P)}$

$det(P)=\begin{vmatrix}
0 & 1  & 2 \\ 
1 & 1  & -1\\ 
1 & -1 & 1 
\end{vmatrix}=-(1+2)+(-1-2)=-6$

$com(P)=\begin{pmatrix}
0 & -2  & -2 \\ 
-3 & -2  & 1\\ 
-3 & 2 & -1 
\end{pmatrix}\Rightarrow t_{com(P)}=\begin{pmatrix}
0 & -3  & -3 \\ 
-2 & -2  & 2\\ 
-2 & 1 & -1 
\end{pmatrix}$

Ainsi $P^{-1}=-\frac{1}{6}\begin{pmatrix}
0 & -3  & -3 \\ 
-2 & -2  & 2\\ 
-2 & 1 & -1 
\end{pmatrix}=\frac{1}{6}\begin{pmatrix}
0 & 3  & 3 \\ 
2 & 2  & -2\\ 
2 & -1 & 1 
\end{pmatrix}$


4. Déduisons $M^{n}$ $\forall n\in \mathbb{N}$.

$M=PDP^{-1}\Rightarrow M^n=PD^nP^{-1}$ avec $D^n=\begin{pmatrix}
0 & 0 & 0 \\ 
0 & 1 & 0 \\ 
0 & 0 & 16^n
\end{pmatrix}$

$M^n=\frac{1}{6}\begin{pmatrix}
0 & 1  & 2 \\ 
1 & 1  & -1\\ 
1 & -1 & 1 
\end{pmatrix}\begin{pmatrix}
0 & 0 & 0 \\ 
0 & 1 & 0 \\ 
0 & 0 & 16^n
\end{pmatrix}\begin{pmatrix}
0 & 3  & 3 \\ 
2 & 2  & -2\\ 
2 & -1 & 1 
\end{pmatrix}$

$M^n=\frac{1}{6}\begin{pmatrix}
2+4.(16)^n & 2-2.(16)^n  & -2+2.(16)^n \\ 
2-2.(16)^n & 2+(16)^n  & -2-(16)^n\\ 
-2+2.(16)^n & -2-(16)^n & 2+(16)^n 
\end{pmatrix}$ $\forall n\in \mathbb{N}$

\medskip
\textbf{Exercice 2}

On considère l'équation différentielle :

\begin{center}
$2y''-y'-y=-7-3x+(12x+5)e^x$.
\end{center}

1. Résolvons l'équation homogène associée à $(E)$. On notera $y_H$ sa solution.

L'équation homogène :
$2y''-y'-y=0$

L'équation caractéristique :

$2r^2-r-1=0$

$\Delta=1+8=9$ $\Longrightarrow r_1=\frac{1-3}{4}=-\frac{1}{2}$ et $r_2=\frac{1+3}{4}=1$.

On a deux racine distinctes donc la solution homogène est de la forme :

$y_H=Ae^{-\frac{1}{2}x}+Be^x$


2. Trouvons une solution particulière de l'équation suivante sous la forme $y_1=mx+p$ avec $m,p\in \mathbb{R}$ :

$2y''-y'-y=-7-3x$.

$y_1=mx+p$

$y'_1=m$

$y''_1=0$.

En reportant dans l'équation précédente, on a :

$0-m-mx-p=-7-3x$.

Par identification :

$\begin{cases}
m=3\\
m+p=7\Longrightarrow p=7-3=4
\end{cases}$

Donc $y_1=3x+4$


3. Trouvons une solution particulière de l'équation suivante sous la forme 

$y_2=x(ax+b)e^x$ avec $a,b\in \mathbb{R}$

$2y''-y'-y=(12x+5)e^x$.

$y_2=x(ax+b)e^x=(ax^2+bx)e^x$

$y'_2=(2ax+b)e^x+(ax^2+bx)e^x=(ax^2+(2a+b)x+b)e^x$

$y''_2=(2ax+2a+b)e^x+(ax^2+(2a+b)x+b)e^x=(ax^2+(4a+b)x+2a+2b)e^x$

En reportant dans l'équation précédente, on obtient :

$2(ax^2+(4a+b)x+2a+2b)e^x-(ax^2+(2a+b)x+b)e^x-(ax^2+bx)e^x=(12x+5)e^x$.

$\Rightarrow (6ax+2a+b)e^x=(12x+5)e^x$

Par identification, on a :

$\begin{cases}
6a=12\\
2a+b=5
\end{cases}\Rightarrow \begin{cases}
a=2\\
b=1
\end{cases}$

Donc $y_2=x(2x+1)e^x$.

4. Déduisons une solution particulière $y_p$ de l'équation $(E)$.

Par superposition des solutions particulières, on obtient :

$y_p=y_1+y_2=3x+4+x(2x+1)e^x$

5. Écrivons la solution générale $y_G$ de l'équation différentielle $(E)$.

$y_G=y_p+y_H=3x+4+x(2x+1)e^x+Ae^{-\frac{1}{2}x}+Be^x$


6. Déterminons la solution $f$ de $(E)$ telle que la tangente à la courbe représentative de $f$ au point $A(0,2)$ soit parallèle à la droite $(D)$ d'équation $y-9x+9=0$.

$y-9x+9=0\Rightarrow y=9x+9$

Au point $A(0,2)$, $f(0)=2\Rightarrow 4 + A +B =2\Rightarrow \underline{A+B=-2}$

La tangente est parallèle à $(D)$ implique que la tangente et $(D)$ ont la même pente qui est 9 selon l'équation de $(D)$.

La pente de la tangente à la courbe de $f$ au point $A$ correspond à $f'(0)$.

D'où $f'(0)=9$

$f(x)=3x+4+x(2x+1)e^x+Ae^{-\frac{1}{2}x}+Be^x$

$f'(x)=3+(4x+1)e^x+(2x^2+x)e^x-\frac{1}{2}Ae^{-\frac{1}{2}x}+Be^x$.

$f'(0)=3+1-\frac{1}{2}A+B=9\Rightarrow \underline{-\frac{1}{2}A+B=5}$

$\begin{cases}
A+B=-2\\
-\frac{1}{2}A+B=5
\end{cases}$

$L_1-L_2\Longrightarrow \frac{3}{2}A=-7\Longrightarrow A=-\frac{14}{3}$

$L_1\Longrightarrow B=-2-A=-2+\frac{14}{3}=\frac{8}{3}$

Ainsi, $f(x)=3x+4+x(2x+1)e^x-\frac{14}{3}e^{-\frac{1}{2}x}+\frac{8}{3}e^x$


\medskip\textbf{Exercice 3}





Soit $f$ une fonction définie sur $[-1, +\infty[$ par :
\[
f(x) = \frac{2}{\pi} \int_{0}^{\frac{\pi}{2}} \sqrt{1 + x \cos^2 t} \, dt.
\]
1. Calcul de $f(-1)$ et $f(0)$

$\cdot$ Pour $x = -1$ :
\[
f(-1) = \frac{2}{\pi} \int_{0}^{\frac{\pi}{2}} \sqrt{1 - \cos^2 t} \, dt.
\]
Or, $\sqrt{1 - \cos^2 t} = \sin t$, donc :
\[
f(-1) = \frac{2}{\pi} \int_{0}^{\frac{\pi}{2}} \sin t \, dt.
\]
La primitive de $\sin t$ est $-\cos t$, d'où :
\[
f(-1) = \frac{2}{\pi} \big[ -\cos t \big]_{0}^{\frac{\pi}{2}} = \frac{2}{\pi} \big( -\cos\frac{\pi}{2} + \cos 0 \big).
\]
En simplifiant, $\cos \frac{\pi}{2} = 0$ et $\cos 0 = 1$, donc :
\[
f(-1) = \frac{2}{\pi} (0 + 1) = \frac{2}{\pi}.
\]

$\cdot$ Pour $x = 0$ :
\[
f(0) = \frac{2}{\pi} \int_{0}^{\frac{\pi}{2}} \sqrt{1 + 0 \cdot \cos^2 t} \, dt = \frac{2}{\pi} \int_{0}^{\frac{\pi}{2}} 1 \, dt.
\]
La primitive de $1$ est $t$, donc :
\[
f(0) = \frac{2}{\pi} \big[ t \big]_0^{\frac{\pi}{2}} = \frac{2}{\pi} \left( \frac{\pi}{2} - 0 \right) = 1.
\]

\textbf{ Conclusion :}
\[
f(-1) = \frac{2}{\pi}, \quad f(0) = 1.
\]

2. Relation entre $f(x)$ et $f\left(\frac{-x}{1+x}\right)$

On cherche à établir la relation :
\[
f(x) = \sqrt{1+x} \cdot f\left(\frac{-x}{1+x}\right).
\]

Définition de $f(x)$ :
\[
f(x) = \frac{2}{\pi} \int_{0}^{\frac{\pi}{2}} \sqrt{1 + x \cos^2 t} \, dt.
\]

Changement de variable :

On pose $u = \arcsin\left(\sqrt{\frac{1}{1+x}} \sin t\right)$, ce qui implique :
\[
\sin t = \sqrt{1+x} \sin u, \quad \cos t = \sqrt{1 - (1+x)\sin^2 u}.
\]
La différentielle devient :
\[
dt = \sqrt{1+x} \cos u \, du.
\]

Les bornes changent comme suit :  
\[
t = 0 \implies u = 0, \quad t = \frac{\pi}{2} \implies u = \arcsin\left(\sqrt{\frac{1}{1+x}}\right) = \frac{\pi}{2}.
\]

Substitution dans l'intégrale :
En remplaçant dans $f(x)$, on obtient :
\[
f(x) = \frac{2}{\pi} \int_{0}^{\frac{\pi}{2}} \sqrt{1+x} \sqrt{1 + \frac{-x}{1+x} \cos^2 u} \cdot \sqrt{1+x} \cos u \, du.
\]

En mettant le facteur $\sqrt{1+x}$ en facteur, on obtient :
\[
f(x) = \sqrt{1+x} \cdot \frac{2}{\pi} \int_{0}^{\frac{\pi}{2}} \sqrt{1 + \frac{-x}{1+x} \cos^2 u} \, du.
\]

En posant $y = \frac{-x}{1+x}$, on reconnaît :
\[
f(x) = \sqrt{1+x} \cdot f\left(\frac{-x}{1+x}\right).
\]

Conclusion :
\[
f(x) = \sqrt{1+x} \cdot f\left(\frac{-x}{1+x}\right).
\]

3. Encadrement pour $-1 \leq x_0 \leq x$ : $0 \leq f(x) - f(x_0) \leq \frac{x - x_0}{4}$

La différence est donnée par :
\[
f(x) - f(x_0) = \frac{2}{\pi} \int_{0}^{\frac{\pi}{2}} \left( \sqrt{1 + x \cos^2 t} - \sqrt{1 + x_0 \cos^2 t} \right) dt.
\]

En utilisant la formule :
\[
\sqrt{a} - \sqrt{b} = \frac{a - b}{\sqrt{a} + \sqrt{b}},
\]
on obtient :
\[
\sqrt{1 + x \cos^2 t} - \sqrt{1 + x_0 \cos^2 t} = \frac{(x - x_0) \cos^2 t}{\sqrt{1 + x \cos^2 t} + \sqrt{1 + x_0 \cos^2 t}}.
\]

Ainsi :
\[
f(x) - f(x_0) = \frac{2}{\pi} \int_{0}^{\frac{\pi}{2}} \frac{(x - x_0) \cos^2 t}{\sqrt{1 + x \cos^2 t} + \sqrt{1 + x_0 \cos^2 t}} \, dt.
\]

Encadrement des termes :
- \( \cos^2 t \leq 1 \), donc le numérateur est majoré par $x - x_0$.
- Le dénominateur est minoré par $2\sqrt{1+x_0}$, avec $1+x_0 \geq \frac{1}{2}$.

En intégrant, on obtient :
\[
f(x) - f(x_0) \leq \frac{x - x_0}{4}.
\]

De plus, \( \sqrt{1 + x \cos^2 t} \geq \sqrt{1 + x_0 \cos^2 t} \), donc \( f(x) - f(x_0) \geq 0 \).

Conclusion :
\[
0 \leq f(x) - f(x_0) \leq \frac{x - x_0}{4}.
\]


4. Encadrement pour $-1 < x < 0$ : $0 \leq f(x) - f(-1) \leq \frac{1}{2}\sqrt{1+x}$

Différence :

\[
f(x) - f(-1) = \frac{2}{\pi} \int_{0}^{\frac{\pi}{2}} \left( \sqrt{1 + x \cos^2 t} - \sqrt{1 - \cos^2 t} \right) dt.
\]

Par des majorations similaires, on montre que :
\[
f(x) - f(-1) \leq \frac{1}{2} \sqrt{1+x}.
\]

5. Continuité de $f$

La fonction sous l'intégrale, $\sqrt{1 + x \cos^2 t}$, est continue en $x$.  
Par le théorème des intégrales dépendant d'un paramètre, $f(x)$ est continue sur $[-1, +\infty[$.

\end{spacing}
\end{document}
